\documentclass[11pt,a4paper]{article}
\usepackage[utf8]{inputenc}
\usepackage[T1]{fontenc}
\usepackage{lmodern}
\usepackage[brazil]{babel}
\usepackage{amsmath,amssymb,amsthm}
\usepackage{geometry}
\usepackage{hyperref}
\usepackage{booktabs}
\usepackage{enumitem}
\usepackage{xcolor}

\geometry{margin=2.5cm}

\definecolor{remark}{rgb}{0.12,0.37,0.55}

\newtheorem{theorem}{Proposição}
\newcommand{\E}{\mathbb{E}}
\newcommand{\dd}{\mathrm{d}}
\newcommand{\hpi}{\hat{\pi}}
\newcommand{\hx}{\hat{x}}

\title{\textbf{Capítulo 6 --- Rigidez Nominal}\\[4pt]
       {\large Derivações e Notas}\\[2pt]
       {\normalsize PPGEco-CCJE-UFES · Macroeconomia I · Romer (2019) Cap.\ 6}}
\author{Projeto de Estudo --- \textit{macro-romer}}
\date{2026}

\begin{document}
\maketitle
\tableofcontents
\bigskip

% ============================================================
\section{Configuração do modelo}
% ============================================================

\subsection{Competição monopolística e demanda Dixit-Stiglitz}

Há um contínuo de firmas no intervalo $[0,1]$, cada uma produzindo uma
variedade diferenciada $Y_i$.  O agregador CES gera demanda para cada
variedade:
\begin{equation}\label{eq:demand}
  Y_i = Y\!\left(\frac{P_i}{P}\right)^{-\eta}, \qquad \eta > 1,
\end{equation}
onde $P = \left[\int_0^1 P_i^{1-\eta}\,di\right]^{1/(1-\eta)}$ é o índice
de preços agregado, $Y$ é a demanda agregada e $\eta$ é a elasticidade-preço
de demanda.

\subsection{Lucro do monopólico}

A firma $i$ maximiza o lucro:
\begin{equation}
  \pi_i = (P_i - MC_i)\,Y_i = \left(P_i - MC_i\right)\!Y\!\left(\frac{P_i}{P}\right)^{-\eta},
\end{equation}
onde $MC_i$ é o custo marginal.  A condição de primeira ordem (preço ótimo
flexível) é:
\begin{equation}\label{eq:optimal_price}
  P_i^* = \frac{\eta}{\eta-1}\,MC_i = \mu\,MC_i, \qquad \mu \equiv \frac{\eta}{\eta-1}.
\end{equation}
O markup $\mu > 1$ é decrescente em $\eta$.  No log:
\begin{equation}
  p_i^* = \log\mu + mc_i \approx \log\mu + p + \text{fatores reais}.
\end{equation}

% ============================================================
\section{Custos de menu (Mankiw 1985)}
% ============================================================

\subsection{Decisão de ajuste}

Suponha que o custo marginal real siga a demanda nominal agregada~$M$
(aproximação): $mc_i \approx m$.  Na presença de um custo fixo de ajuste
$z > 0$ (custo de menu), a firma compara o \emph{ganho de lucro} de
re-otimizar contra $z$.

Expandindo o lucro em segunda ordem ao redor de $P_i^*$:
\begin{equation}\label{eq:profit_gain}
  \Pi(P_i) \approx \Pi(P_i^*) - \frac{\eta-1}{2}\,(p_i - p_i^*)^2,
\end{equation}
pois a derivada de segunda ordem do lucro em relação ao log-preço é
$-(\eta-1)\Pi^*/P_i^{*2}$, normalizada a $-(\eta-1)$.  Portanto, o ganho
de reajustar é:
\begin{equation}\label{eq:G}
  G(p_{dev}) \equiv \Pi(P_i^*) - \Pi(P_i) = \frac{\eta-1}{2}\,p_{dev}^2,
\end{equation}
onde $p_{dev} = p_i - p_i^*$ é o log-desvio do preço atual em relação ao
ótimo.

\subsection{Limiar de ajuste}

A firma ajusta se e somente se $G(p_{dev}) \geq z$:
\begin{equation}\label{eq:threshold}
  \frac{\eta-1}{2}\,p_{dev}^2 \geq z
  \quad\Longleftrightarrow\quad
  |p_{dev}| \geq \bar{p}(z) \equiv \sqrt{\frac{2z}{\eta-1}}.
\end{equation}
Observe que $\bar{p}(z)$ cresce com $z$ (mais rigidez para custo maior) e
decresce com $\eta$ (firmas com demanda mais elástica ajustam com mais
frequência).

\subsection{Externalidade de demanda agregada (Ball-Romer)}

Com $n$ firmas simétricas, cada firma ignora o efeito do seu preço sobre a
demanda das demais.  O ganho \emph{social} de ajustar é:
\begin{equation}
  G^{social} = n\,G^{private} = n\,\frac{\eta-1}{2}\,p_{dev}^2.
\end{equation}

\textbf{Falha de coordenação:} existe uma região em que
$G^{private} < z \leq G^{social}$.  Nessa região nenhuma firma ajusta
individualmente, mesmo que todos se beneficiassem se ajustassem
simultaneamente.  Isso gera \textbf{rigidez de preços agregada} mesmo quando
os custos de menu individuais são pequenos --- o resultado central de Mankiw
(1985) e Ball-Romer (1990).

% ============================================================
\section{Precificação de Calvo}
% ============================================================

\subsection{Estrutura}

Em cada período uma firma mantém seu preço com probabilidade $\alpha$ (o
processo de Calvo é independente de estado e do tempo desde o último ajuste).
Com probabilidade $(1-\alpha)$ a firma re-otimiza escolhendo $\tilde{P}_t$.

A probabilidade de que um preço fixado $s$ períodos atrás ainda vigore é:
\begin{equation}
  P(\text{preço inalterado por }s\text{ períodos}) = \alpha^s.
\end{equation}

\subsection{Preço de re-otimização}

A firma que re-otimiza em $t$ escolhe $\tilde{P}_t$ para maximizar o valor
presente esperado dos lucros futuros, ponderado pela probabilidade de
permanência:
\begin{equation}\label{eq:reset_price}
  \log\tilde{P}_t = (1-\alpha\beta)\sum_{s=0}^{\infty}(\alpha\beta)^s
                    \E_t\!\left[\log P_{t+s}^* \right],
\end{equation}
onde $P_{t+s}^* = \mu\,MC_{t+s}$ é o preço ótimo flexível no período $t+s$.

\subsection{Índice de preços agregado}

O índice CES agrega preços antigos (peso $\alpha$) e novos (peso $1-\alpha$):
\begin{equation}\label{eq:price_index}
  P_t^{1-\eta} = \alpha\,P_{t-1}^{1-\eta} + (1-\alpha)\,\tilde{P}_t^{1-\eta}.
\end{equation}
Log-linearizando (em torno de inflação zero):
\begin{equation}\label{eq:price_index_log}
  p_t = \alpha\,p_{t-1} + (1-\alpha)\,\tilde{p}_t.
\end{equation}

% ============================================================
\section{Curva de Phillips Novo-Keynesiana}
% ============================================================

\subsection{Derivação da CPNK}

Definindo inflação $\pi_t = p_t - p_{t-1}$ e combinando
\eqref{eq:reset_price} com \eqref{eq:price_index_log}:

\begin{equation}
  \tilde{p}_t - p_{t-1} = (1-\alpha\beta)\sum_{s=0}^{\infty}(\alpha\beta)^s
                           \E_t[\pi_{t+1:t+s} + \hat{mc}_{t+s}],
\end{equation}
onde $\hat{mc}_t$ é o custo marginal em log-desvio do estado estacionário.

Após manipulação algébrica obtém-se a forma recursiva:
\begin{equation}\label{eq:nkpc}
  \boxed{\pi_t = \beta\,\E_t[\pi_{t+1}] + \kappa\,x_t + u_t},
\end{equation}
onde $x_t$ é o hiato do produto e a inclinação é:
\begin{equation}\label{eq:kappa}
  \kappa = \frac{(1-\alpha)(1-\alpha\beta)}{\alpha}\,(\omega + \sigma^{-1}).
\end{equation}
Aqui $\sigma^{-1}$ é a elasticidade de substituição intertemporal e $\omega$
é a inversa da elasticidade de Frisch (do trabalho).

\subsection{Interpretação dos parâmetros}

\begin{itemize}[leftmargin=*]
  \item $\alpha \nearrow$: mais rigidez → CPNK mais plana ($\kappa \searrow$).
  \item $\beta \nearrow$: firmas pesam mais o futuro → expectativas mais
        importantes.
  \item $\omega \nearrow$ ou $\sigma \searrow$: custo marginal sobe mais com
        produto → CPNK mais íngreme.
\end{itemize}

\begin{table}[h!]
\centering
\caption{Parâmetros da CPNK: valores de referência}
\begin{tabular}{lll}
\toprule
Parâmetro & Valor & Interpretação \\
\midrule
$\alpha$  & 0,75  & 75\,\% das firmas não ajustam em cada trimestre \\
$\beta$   & 0,99  & Fator de desconto trimestral \\
$\sigma$  & 1,0   & IES = 1 (log-utilidade) \\
$\omega$  & 1,0   & Elasticidade de Frisch inversa \\
$\kappa$  & $\approx 0,085$ & Inclinação derivada \\
\bottomrule
\end{tabular}
\end{table}

% ============================================================
\section{Rigidezes reais e falha de coordenação (Ball-Romer)}
% ============================================================

\subsection{Rigidezes nominais vs.\ reais}

Mankiw (1985) mostra que rigidezes nominais (custos de menu) geram
não-neutralidade da moeda mesmo quando os custos individuais são pequenos,
porque a externalidade de demanda agrega as perdas sociais.

Ball e Romer (1990) mostram que \textbf{rigidezes reais} (e.g.\ rigidez real
de salários, competição imperfeita) amplificam o efeito de custos de menu
pequenos.  Com rigidezes reais, o preço ótimo flexível $P^*$ responde pouco
ao ambiente macroeconômico, de modo que o custo de não ajustar é pequeno e
mesmo custos de menu ínfimos geram rigidez agregada substancial.

\subsection{Condição de falha de coordenação}

Definindo o efeito estratégico $\gamma \equiv \partial P_i^*/\partial P$
(resposta do preço ótimo individual ao nível de preços agregado), a condição
necessária para falha de coordenação é:
\begin{equation}
  \gamma > 0 \quad \text{e} \quad z < G^{social} - G^{private} = (n-1)\,G^{private}.
\end{equation}
Quando $\gamma \in (0,1)$ (complementaridade estratégica parcial), o
equilíbrio com preços fixos pode ser auto-realizável, dando origem a
múltiplos equilíbrios e recessões de equilíbrio.

% ============================================================
\section{Quadro resumo e calibração brasileira}
% ============================================================

\begin{table}[h!]
\centering
\caption{Calibração para o Brasil (dados BCB)}
\begin{tabular}{llll}
\toprule
Série & Código BCB & Período & Uso \\
\midrule
IPCA & SGS 433 & 2000--2024 & Inflação realizada \\
Selic & SGS 4189 & 2000--2024 & Taxa nominal de referência \\
\bottomrule
\end{tabular}
\end{table}

A CPNK calibrada com $\alpha=0{,}75$ e parâmetros padrão implica:
\begin{equation}
  \kappa = \frac{(1-0{,}75)(1-0{,}75 \times 0{,}99)}{0{,}75}(1 + 1)
         \approx 0{,}085.
\end{equation}

Regra de Taylor simplificada para comparação:
\begin{equation}
  i_t = \bar{r} + \phi_\pi\,\pi_t, \qquad
  \bar{r} = 0{,}75\%\text{ a.m.}, \quad \phi_\pi = 1{,}5.
\end{equation}

\end{document}
